% Bloque C: Nicolás, 4:35. S9--S12. RelayBook: relays, state-based resume,
% réplicas por región, hipótesis y métricas.
\section{La propuesta}
\begin{frame}{RelayBook: recuperar en relays regionales}
  \centering% S9. Binance hoy vs RelayBook: dónde termina la recuperación.
\begin{tikzpicture}[x=1cm,y=1cm]
  \path[use as bounding box] (0,-0.5) rectangle (13.5,4.65);
  % ---------------- (a) Binance hoy ----------------
  \node[anchor=west,font=\small\bfseries] at (0,4.45) {(a) Binance hoy};
  \node[malo,font=\scriptsize,text width=4.2cm] (o) at (3.1,3.55) {\textbf{Servidor central}\\stream a todos + snapshots REST};
  \foreach \i/\x in {1/0.55,2/1.8,3/3.05,4/4.3,5/5.55} {
    \node[cliente] (c\i) at (\x,0.75) {cliente};
    \draw[flow,colAcento,thin] (\x-0.12,3.05) -- (c\i.north west);
    \draw[rec,colWake] (c\i.north east) -- (\x+0.12,3.05);
  }
  \node[font=\scriptsize,text=colWake] at (3.05,0.1) {cada reconexión pide el libro al servidor central};
  \draw[colNeutro!40] (6.55,-0.25) -- (6.55,4.6);
  % ---------------- (b) RelayBook ----------------
  \node[anchor=west,font=\small\bfseries,text=colAcento] at (6.85,4.45) {(b) RelayBook};
  \node[box,font=\scriptsize,text width=4.2cm] (or) at (10.2,3.75) {\textbf{Origen}\\un solo feed con secuencia};
  \node[acento,font=\scriptsize,text width=2.5cm] (ra) at (8.45,2.2) {\textbf{Relay} región A\\copia del libro};
  \node[acento,font=\scriptsize,text width=2.5cm] (rb) at (11.95,2.2) {\textbf{Relay} región B\\copia del libro};
  \draw[flow,colAcento] (or) -- (ra);
  \draw[flow,colAcento] (or) -- (rb);
  \foreach \r/\xs in {ra/7.45,rb/10.95} {
    \foreach \k in {0,1,2} {
      \pgfmathsetmacro{\xc}{\xs+\k*1.0}
      \node[cliente,minimum width=0.85cm] (k) at (\xc,0.75) {cl.};
      \draw[flow,colAcento,thin] (\xc-0.1,1.65) -- (\xc-0.1,1.03);
      \draw[rec,colReq] (\xc+0.12,1.03) -- (\xc+0.12,1.65);
    }
  }
  \node[font=\scriptsize,text=colReq] at (10.2,0.1) {la reconexión se resuelve en el relay};
  % Leyenda
  \node[anchor=west,font=\scriptsize] at (0.2,-0.3) {%
    \tikz\draw[flow,colAcento] (0,0) -- (0.5,0); distribución \quad
    \tikz\draw[rec,colWake] (0,0) -- (0.5,0); recuperación en el servidor central \quad
    \tikz\draw[rec,colReq] (0,0) -- (0.5,0); recuperación en el relay};
\end{tikzpicture}

  \takeaway{Idea 1: separar la \emph{distribución} de la \emph{recuperación}.}
  \note{Nicolás. 7:05--8:10 (65 s). A la izquierda, Binance hoy: el mismo
  servidor central reparte el stream a todos los clientes y además atiende
  cada recuperación; son las flechas rojas. A la derecha, RelayBook. El
  origen publica un solo feed con números de secuencia hacia una capa de
  relays, uno o más por región. Cada relay mantiene una copia exacta del
  libro y atiende a sus clientes. Cuando un cliente se reconecta, lo
  resuelve el relay, en verde, y el origen ni se entera. Así el trabajo
  del origen depende del número de relays, no del número de clientes ni
  de cuántos se desconectan. El costo: una capa nueva de servidores que
  hay que operar, y que puede fallar; por eso un cliente puede reanudar en
  cualquier otro relay.}
\end{frame}

\begin{frame}{Idea 2: \emph{state-based resume}, reanudar por estado}
  \centering% S11. State-based resume en 4 pasos, con el mismo libro de la S3--S5.
\begin{tikzpicture}[x=1cm,y=1cm]
  \path[use as bounding box] (0,0.1) rectangle (13.5,5.5);
  \def\pv#1{\textcolor{colVenta}{#1}}
  \def\pc#1{\textcolor{colCompra}{#1}}
  \def\ok#1{\textcolor{colReq}{\textbf{#1}}}
  \tikzset{
    paso/.style={font=\scriptsize,align=left,anchor=north west,inner sep=0pt},
    pie/.style={font=\scriptsize,align=left,anchor=north west,inner sep=0pt},
    panel/.style={box,font=\scriptsize,anchor=north,inner sep=3pt},
  }
  % ---------- 1. Antes del corte ----------
  \node[paso,text width=2.45cm] at (0.05,5.5) {\textbf{1.} Antes del corte, el cliente tiene el libro al día hasta el \#1000.};
  \node[panel,acento,minimum width=2.5cm] (p1) at (1.3,4.3) {\textbf{Copia del cliente}\\[1pt]
    \begin{tabular}{@{}l@{\quad}r@{}}\pv{67\,002}&0.5\\\pv{67\,001}&2.0\\\pc{67\,000}&1.2\\\pc{66\,999}&0.8\end{tabular}};
  % ---------- 2. Corte ----------
  \node[paso,text width=2.95cm] at (2.9,5.5) {\textbf{2.} El cliente se desconecta y no recibe 6 mensajes.};
  \node[panel,malo,dashed,minimum width=3.0cm] (p2) at (4.4,4.3) {\begin{tabular}{@{}l@{\ \ }l@{}}\#1001&67\,000 $\to$ 1.0\\\#1002&67\,000 $\to$ 1.5\\\#1003&67\,000 $\to$ 0.7\\\#1004&67\,000 $\to$ 0.9\\\#1005&66\,999 $\to$ 0.6\\\#1006&67\,002 $\to$ 0\end{tabular}};
  \node[pie,text width=3.0cm,text=colWake] at (2.9,2.0) {67\,000 cambió 4 veces\\y 67\,002 se borró.};
  % ---------- 3. Tabla del relay ----------
  \node[paso,text width=3.65cm] at (6.25,5.5) {\textbf{3.} El relay recibe los 6 mensajes y guarda por nivel su valor actual y el \textbf{último mensaje} que lo cambió.};
  \node[panel,minimum width=3.7cm] (p3) at (8.1,4.3) {\textbf{Tabla del relay}\\[1pt]
    \begin{tabular}{@{}l@{\ \ }c@{\ \ }r@{}}
      Precio & Cantidad & Último\\\hline
      \pv{67\,002} & 0 \emph{tombstone} & \ok{\#1006}\\
      \pv{67\,001} & 2.0 & \#980\\
      \pc{67\,000} & 0.9 & \ok{\#1004}\\
      \pc{66\,999} & 0.6 & \ok{\#1005}
    \end{tabular}};
  \node[pie,text width=3.7cm] at (6.25,2.2) {\textcolor{colReq}{En verde}, cambios después del \#1000.\\Un \textbf{tombstone} es un nivel borrado guardado en 0.};
  % ---------- 4. Respuesta ----------
  \node[paso,text width=3.15cm] at (10.3,5.5) {\textbf{4.} El cliente vuelve con el \#1000 y el relay envía solo lo que cambió después.};
  \node[panel,bueno,minimum width=3.15cm] (p4) at (11.875,4.3) {\textbf{Respuesta con 3 niveles}\\[1pt]
    \begin{tabular}{@{}l@{\ \ }l@{}}\pv{67\,002}&$\to$ borrar\\\pc{67\,000}&$\to$ 0.9\\\pc{66\,999}&$\to$ 0.6\end{tabular}};
  \node[pie,text width=3.15cm,text=colReq] at (10.3,2.75) {1.0, 1.5 y 0.7 no viajan porque solo cuenta el último valor.\\El cliente queda igual\\al relay \okmark};
  % ---------- Flechas ----------
  \draw[flow,colWake,dashed] (2.55,3.45) -- (2.9,3.45);
  \draw[flow] (5.9,3.45) -- (6.25,3.45);
  \draw[flow,colReq] (9.95,3.45) -- (10.3,3.45);
  % ---------- Comparación ----------
  \node[font=\scriptsize\bfseries,align=left,text width=2.2cm,anchor=west] at (0.05,0.5) {Recuperación del mismo corte};
  \node[malo,font=\scriptsize,text width=3.2cm] at (4.25,0.5) {\textbf{Binance} envía el libro completo, 2\,000 niveles};
  \node[box,font=\scriptsize,text width=3.2cm] at (7.95,0.5) {\textbf{Estilo Ably} envía los 6 mensajes perdidos};
  \node[bueno,font=\scriptsize,text width=3.2cm] at (11.65,0.5) {\textbf{RelayBook} envía solo 3 niveles};
\end{tikzpicture}

  \takeaway{Se manda el valor actual de cada nivel que cambió, no la historia.}
  \source{Si los cambios pesan tanto como el libro, el relay manda el snapshot completo: reanudar nunca cuesta más que resincronizar.}
  \note{Nicolás. 8:10--9:40 (90 s). Esta es la idea central. El relay guarda,
  además del valor de cada nivel, en qué mensaje cambió por última vez. Un
  cliente se desconecta después del mensaje 1000. Mientras no está, el
  nivel 67 000 cambia cuatro veces, el 66 999 una vez, y el 67 002
  desaparece: el relay no lo borra, lo deja en cero, una lápida, para poder
  avisar. El cliente vuelve y dice: lo último que vi fue el 1000. El relay
  se queda con los niveles cuyo último cambio es mayor que 1000: tres
  niveles. Le manda su valor actual y el cliente queda idéntico al relay.
  Funciona porque cada delta trae el valor nuevo, no una resta: de 1.0,
  1.5, 0.7 y 0.9 solo importa el 0.9. Comparen abajo: Binance manda el
  libro entero, un resume estilo Ably manda los seis mensajes perdidos, y
  nosotros tres niveles. Y una regla de seguridad: si la lista ya pesa como
  el libro, el relay manda el libro, pero desde el relay, nunca desde el
  origen.}
\end{frame}

\begin{frame}{Idea 3: réplicas por región (\emph{leader-follower})}
  \centering% S11. Réplicas por región: leader-follower, lag acotado, gaps, checksum, conflation.
\begin{tikzpicture}[x=1cm,y=1cm]
  \path[use as bounding box] (0,0.05) rectangle (13.5,4.6);
  \tikzset{rl/.style={font=\scriptsize,text width=1.2cm,inner sep=2pt}}
  % ---------------- Diagrama ----------------
  \node[acento,font=\scriptsize,text width=3.6cm] (or) at (3.3,4.05) {\textbf{Origen = líder}\\numera cada mensaje};
  \draw[colNeutro!60,rounded corners=3pt] (0.05,0.1) rectangle (3.2,2.95);
  \draw[colNeutro!60,rounded corners=3pt] (3.45,0.1) rectangle (6.6,2.95);
  \node[font=\scriptsize\bfseries] at (1.62,2.72) {Región A};
  \node[font=\scriptsize\bfseries] at (5.02,2.72) {Región B};
  \node[bueno,rl] (a1) at (0.85,1.85) {Relay 1\\sano};
  \node[malo,rl]  (a2) at (2.4,1.85) {Relay 2\\lag $>$ 500\,ms};
  \node[bueno,rl] (b1) at (4.25,1.85) {Relay 3\\sano};
  \node[bueno,rl] (b2) at (5.8,1.85) {Relay 4\\sano};
  \foreach \r in {a1,a2,b1,b2} \draw[flow,colAcento] (or.south) -- (\r.north);
  \node[font=\tiny,text=colAcento,fill=white,inner sep=1pt] at (3.3,3.25) {feed = log de replicación};
  \node[cliente] (cl) at (2.4,0.55) {cliente};
  \draw[colWake,dashed] (a2.south) -- (cl.north);
  \draw[rec,colReq] (cl.west) to[out=180,in=270] node[pos=0.35,below,font=\tiny,text=colReq] {reanuda aquí} (a1.south);
  \foreach \x in {4.25,5.8} {
    \node[cliente,minimum width=0.8cm] at (\x,0.55) {cl.};
    \draw[flow,colAcento,thin] (\x,1.3) -- (\x,0.83);
  }
  % ---------------- Mecanismos ----------------
  \tikzset{mec/.style={box,font=\scriptsize,text width=6.2cm,align=left,anchor=west}}
  \node[mec] at (6.95,4.15) {\textbf{1. Un solo orden:} solo el origen numera; todos aplican lo mismo.};
  \node[mec] at (6.95,3.25) {\textbf{2. Lag acotado:} si un relay se atrasa $>$\,500\,ms, sus clientes van a otro relay.};
  \node[mec] at (6.95,2.35) {\textbf{3. Gaps:} redirige o espera el snapshot del feed (cada 1\,s); nunca pide al origen.};
  \node[mec] at (6.95,1.45) {\textbf{4. Checksum} de los 10 mejores niveles: detecta copias que divergen.};
  \node[mec,bueno,align=left] at (6.95,0.55) {\textbf{5. Conflation:} al cliente lento no se le corta; recibe solo el último valor.};
\end{tikzpicture}

  \takeaway{Un relay puede ir atrasado, pero nunca muestra un libro que no existió.}
  \note{Nicolás. 9:40--10:50 (70 s). Los relays son réplicas, como en una base
  de datos replicada: el origen es el líder, los relays son seguidores, y
  el feed con secuencia es el log de replicación. Como solo el origen
  numera los mensajes, todos aplican lo mismo en el mismo orden: un relay
  puede ir unos milisegundos atrasado, pero nunca inventa un libro. Ese
  atraso se llama replication lag, y lo acotamos: si pasa de 500 ms, la
  mitad de nuestra meta de un segundo, el relay se marca no saludable y
  sus clientes reanudan en otro relay de la misma región. Si un relay
  pierde mensajes del feed, no le pide nada al origen: redirige a sus
  clientes o espera el snapshot que el feed manda cada segundo. Un checksum
  detecta copias que divergen. Y a los clientes lentos no se les corta:
  el relay junta los cambios y les manda solo el último valor de cada
  nivel, lo que llamamos conflation.}
\end{frame}

\begin{frame}{Hipótesis y cómo la medimos}
  \small
  \textbf{Parte 1.} El contrato de Binance, con tráfico real acelerado 10$\times$,
  \textbf{sigue degradado} cuando termina el pico.\par
  \textbf{Parte 2.} En las mismas condiciones, RelayBook \textbf{cumple las tres
  metas}, y un resume estilo Ably tiene \textbf{más atraso}, porque cada
  recuperación arrastra todos los mensajes perdidos.\par
  \medskip
  \centering
  \begin{tabular}{p{3.2cm}p{7.6cm}l}
    \toprule
    Métrica & Qué mide & Meta \\
    \midrule
    \emph{Full resync ratio} & De las reconexiones durante el pico, cuántas terminan en un snapshot completo & $< 1\%$ \\
    \emph{Time to recover} & Desde el fin del pico, cuánto tarda el 99\% de los clientes en tener un libro válido & $\leq 30$\,s \\
    \emph{Staleness} p99 & Atraso del libro del cliente frente al origen; el 99\% de las muestras está por debajo & $< 1$\,s \\
    \bottomrule
  \end{tabular}
  \takeaway{Comparamos tres diseños: Binance, estilo Ably y RelayBook.}
  \note{Nicolás. 10:50--11:40 (50 s). La hipótesis tiene dos partes, y
  cualquiera puede salir refutada. Primera: el contrato de Binance, con
  tráfico real acelerado diez veces, sigue degradado después del pico.
  Segunda: RelayBook cumple las tres metas, y el resume estilo Ably tiene
  más atraso. Las métricas son tres. Full resync: qué fracción de las
  reconexiones termina en un libro completo; queremos menos del uno por
  ciento. Time to recover: cuánto tarda el 99\% de los clientes en volver a
  tener un libro válido; queremos 30 segundos o menos, sin intervención.
  Staleness p99: cuánto atraso tiene el libro del cliente; queremos menos
  de un segundo. Ably es la referencia porque reanuda reenviando los
  mensajes perdidos dentro de una ventana de dos minutos. Pedro explica
  cómo lo probamos.}
\end{frame}
